\documentclass[10pt,a4paper]{amsart}
\usepackage[margin=19mm]{geometry}
\usepackage{amsmath,amssymb,mathtools}
\usepackage[scr=rsfs,cal=euler]{mathalfa}
\usepackage{xfrac}
\usepackage[unicode,hidelinks,bookmarks=false]{hyperref}
\newcommand{\ie}{i.e.\ }
\newcommand{\cf}{cf.\ }
\newcommand{\resp}{respectively\ }
\newcommand{\Subsection}{\subsection}
\providecommand{\nobreakdash}{\nobreak\hbox{-}}
\setlength{\emergencystretch}{2em}
\multlinegap0pt
\usepackage{bm}
\usepackage{bbm}
\usepackage{dsfont}
\usepackage{xspace}
\usepackage{cancel}
\usepackage{mathtools}
\usepackage[shortlabels]{enumitem}
\usepackage{amsthm}
\makeatletter
\@ifundefined{theorem}{\newtheorem{theorem}{Theorem}}{}
\@ifundefined{lemma}{\newtheorem{lemma}[theorem]{Lemma}}{}
\makeatother
\newcommand{\B}[1]{{\mathbf #1}}
\newcommand{\C}[1]{{\mathscr #1}}
\title[Torsion in asymptotic cones of free groups]{Torsion in asymptotic cones of free groups}
\author{Jarek K\k{e}dra, Assaf Libman, Zipei Nie}
\date{Source: arXiv:2608.10854v2 (CC BY 4.0)}
\begin{document}
\maketitle

\noindent\emph{Source and licence.} Torsion in asymptotic cones of free groups. Version arXiv:2608.10854v2 (CC BY 4.0), distributed under the Creative Commons Attribution 4.0 International licence, as recorded in the arXiv licence metadata for this exact version and shown on the abstract page https://arxiv.org/abs/2608.10854v2. Full rights notice: licenses/arxiv-2608.10854-CC-BY-4.0.txt.\par\smallskip

\noindent\textit{Editorial scope.} Counted unit: the Lemma whose source label is \texttt{L:intricate} in the article (source0.tex lines 1348--1366, source label \texttt{L:intricate}), delivered together with the article's own complete original proof of it (source0.tex lines 1367--1564, one \texttt{proof} environment of 198 lines). No mathematics is written, repaired or re-derived: statement and proof are the article's own text line for line. Required context retained uncounted: eqn:phi\_k(T\_j) (lines 899--903); L:fi-pi (lines 1011--1022); L:fi-pi:2 (lines 1013--1021); L:fik-anti-triangle (lines 1252--1257). Every definition, display and label of those lines is retained unchanged. Omitted from this package and not redistributed: 1--898, 904--1010, 1022--1251, 1258--1347, 1565--1885. Nothing inside a retained excerpt is omitted or altered: every retained body is delivered exactly as the article's lines are written, so no omission receipt of any kind accompanies this package. Citations: the retained excerpts carry no citation command at all, so no bibliography entry of the article is transcribed and no reference is redistributed. Reference adaptation: none was needed and none was made; every reference inside the delivered unit resolves inside it, so no unresolved reference or citation is produced. Printed slips kept literal: no printed word, symbol, hypothesis or number of the counted statement or of its proof was altered. Layout and class: the article's own class is \texttt{article} with packages including amsmath, amsthm, amssymb, authblk, color, soul, xcolor, enumitem, euscript, fontenc, graphics, fullpage, hyperref, mathtools; the wrapper is the lane's pinned \texttt{amsart} editorial wrapper at 10pt on A4, which restates the theorem-like environments the retained text needs and adds the article's own macro definitions that the retained text needs (\texttt{\textbackslash B}, \texttt{\textbackslash C}), with extra wrapper packages (bm, bbm, dsfont, xspace, cancel, mathtools, enumitem). The delivered numbering is the unit's own. Assets: no retained span contains a figure environment, a graphics-inclusion command, a PDF-page inclusion, a table or a TikZ picture; no image file is copied into this package, the package declares no asset dependency and no image file is redistributed with it. Licence: version arXiv:2608.10854v2 of this article is distributed under the Creative Commons Attribution 4.0 International licence, as recorded in the arXiv licence metadata for this exact version and shown on the abstract page https://arxiv.org/abs/2608.10854v2. A full rights notice is delivered at licenses/arxiv-2608.10854-CC-BY-4.0.txt. Characters: every retained excerpt is pure ASCII, so the delivered engine, XeLaTeX with its default Latin Modern text font, reports no missing character.
\begin{equation}\label{eqn:phi_k(T_j)}
\rho_k(T_j) =
\rho_k((x_jy_j)^n (y_jx_j)^{-n}) =
(x_{j \, {\rm mod} \, 2})^{2n} (x_{j \, {\rm mod} \, 2})^{-2n} =1.
\end{equation}

\begin{lemma}\label{L:fi-pi}
Let $2\leq k\leq 2n$.
\begin{enumerate}[label=(\arabic*)]
\item \label{L:fi-pi:1}
$\pi_I\circ \varphi^{(k)}=\pi_I$ for any $I\subseteq\{0,1,\ldots,k\}$.
\item \label{L:fi-pi:2}
$\widehat{\pi}_{\{j\}}\circ\varphi^{(k)}\circ\varphi_j(g)=\varphi^{(k)}(g)$ for
any $1\leq j\leq k$ and any $g \in \B F_{\mathbb N\setminus\{j\}}$.
\item \label{L:fi-pi:3}
$\rho_k\circ\varphi^{(k)} = \varphi^{(k)}\circ \rho_k$.
\end{enumerate}
\end{lemma}

\begin{lemma}\label{L:fik-anti-triangle}
With the notation and assumptions above we have
$$
\|\varphi^{(k)}(w)\|\geq \|\varphi^{(k)}(\rho_{\infty}(w_0))\| + \|\varphi^{(k)}(\rho_{\infty}(w_Lw_R))\|.
$$
\end{lemma}

\begin{lemma}\label{L:intricate}
Let $3\leq k\leq 2n-1$ be an odd integer. Let $w\in H=\langle x_0,x_1\rangle$ be a non-trivial
reduced word which starts and ends with the same letter, including sign.
Then either
$$
\|\varphi^{(k-1)}(w)\|\geq 4n + \|\varphi^{(k+1)}(w)\|
$$
or there exist $w_0,w_L,w_R\in H$ such that 
\begin{enumerate}[label=(\alph*)]
\item $w=w_Lw_0w_R$; equality in $H$, not necessarily a decomposition of $w$ into the concatenation of three subwords. % (Statement rephrased in response to Nie's request to make it clearer. I hope it is clearer now.)
% \AL{The concatenation $w_Lw_0w_R$ need not be a reduced word in the alphabet $x_0,x_1$ and hence not equal to $w$ as words, but its reduced form is $w$.}
% not necessarily as concatenation of words, 

\item $w_0\neq 1$ and $w_Lw_R\neq 1$, and 
$$
\|\varphi^{(k-1)}(w)\|\geq  \|\varphi^{(k+1)}(w_0)\|+ \|\varphi^{(k+1)}(w_Lw_R)\|.
$$
\end{enumerate}
\end{lemma}

\begin{proof}
Let $w=t_1\dots t_m$, where $t_i\in \{x_0^{\pm 1},x_1^{\pm 1}\}$. 
After replacing $w$ with
$w^{-1}$, we can assume, without loss of generality, that $w$ starts and ends 
with either $x_0$ or $x_1$. The proof has
two parts depending on whether $t_1$ is equal to $x_0$ or $x_1$. 
We will conduct both cases in parallel as follows.

Since the homomorphisms
$\widehat{\pi}_{\{k\}}$ and $\widehat{\pi}_{\{k+1\}}$ 
map generators to generators, they are Lipschitz with
constant $1$, which is the first inequality in the following
parallel computation; The second equalities follow from Lemma \ref{L:fi-pi}\ref{L:fi-pi:2}
$$
\begin{array}{cc}
\parbox[t]{0.45 \textwidth}{
\begin{align*}
t_1 &=x_0\\
\hline\\
\|\varphi^{(k-1)}(w)\|
&\geq \|\widehat{\pi}_{\{k+1\}}\varphi^{(k-1)}(w)\|\\
&=\|\widehat{\pi}_{\{k+1\}}\varphi^{(k+1)}\varphi_{k+1}\varphi_k(w)\|\\
&=\|\varphi^{(k+1)}\varphi_k(w)\|.
\end{align*}
}&
\parbox[t]{0.45 \textwidth}{
\begin{align*}
t_1 &= x_1\\
\hline\\
\|\varphi^{(k-1)}(w)\|
&\geq \|\widehat{\pi}_{\{k\}}\varphi^{(k-1)}(w)\|\\
&=\|\widehat{\pi}_{\{k\}}\varphi^{(k)}\varphi_k(w)\|\\
&=\|\varphi^{(k)}(w)\|\\
&=\|\varphi^{(k+1)}\varphi_{k+1}(w)\|.
\end{align*}
}
\end{array}
$$
Let $j = k$ if $t_1=x_0$ and $j = k+1$ if $t_1=x_1$.
Recall that $k$ is odd, so 
\[
t_1 = t_m=x_{(j+1)\,{\rm mod}\,2}.
\]
In this notation both estimates above can be written together as
\begin{equation}
\|\varphi^{(k-1)}(w)\|\geq \|\varphi^{(k+1)}\varphi_{j}(w)\|
\label{Eq:j-notation}
\end{equation}
and we will now analyse the element $\varphi^{(k+1)}\varphi_{j}(w)$.

Let $(c,\C F)$ be a cancellation system for
$\varphi^{(k+1)}\varphi_{j}(w)$ 
with minimal $c$.
Then $\pi_{\{j\}}(c)$ and $\widehat{\pi}_{\{j\}}(c)$ are cancellation 
sequences for
$\pi_{\{j\}}(\varphi^{(k+1)}\varphi_{j}(w))$ and 
$\widehat{\pi}_{\{j\}}(\varphi^{(k+1)}\varphi_{j}(w))$, respectively.
Moreover, since $j=k$ or $j=k+1$, by Lemma \ref{L:fi-pi}\ref{L:fi-pi:2}
$$
\widehat{\pi}_{\{j\}}(\varphi^{(k+1)}\varphi_{j}(w)) 
=\varphi^{(k+1)}(w).
$$
It follows that
\[
\ell(\widehat{\pi}_{\{j\}}(c))\geq \|\varphi^{(k+1)}(w)\|.
\]
%
%
\subsubsection*{Case 1: $\ell(\pi_{\{j\}}(c))\geq 4n$.}
%Assume first that $\ell(\pi_{\{j\}}(c))\geq 4n$. 
Then
\begin{align*}
\|\varphi^{(k-1)}(w)\|
&\geq\|\varphi^{(k+1)}\varphi_{j}(w)\| &\text{by \eqref{Eq:j-notation}} \\
&= \ell(c) = \ell(\pi_{\{j\}}(c)) + \ell(\widehat{\pi}_{\{j\}}(c))\\ 
&\geq 4n + \|\varphi^{(k+1)}(w)\|,
\end{align*}
which proves the first alternative in the statement of the lemma.


\subsubsection*{Case 2: {$\ell(\pi_{\{j\}}(c))< 4n$.}}
Recall that if $t_1=x_0$ then $j$ is odd and if $t_1=x_1$ then $j$ is even.
Since $\varphi_j(t_i) = t_iT_j$ or $T_j^{-1}t_i$, it follows that in both cases
$$
\varphi_j(w) = \varphi_j(t_1\cdots t_m) = t_1T'_1t_2T'_2 \cdots t_mT'_m,
$$
where $T'_i = T_j^{\pm 1}$ or $T'_i=1$, and $T'_m=T_j$.
Consequently,
$$
\varphi^{(k+1)}(\varphi_j(w)) = \varphi^{(k+1)}(t_1)\varphi^{(k+1)}(T'_1) \cdots 
\varphi^{(k+1)}(t_m)\varphi^{(k+1)}(T'_m).
$$
Since $T'_m=T_j$ is a subsequence in $\varphi^{(k+1)}(T'_m)$ and since 
$T_j=(x_jy_j)^n(x_j^{-1}y_j^{-1})^n$
is a reduced  word of length $4n$, the cancellation sequence $c$ must leave at least
one letter $x_j^{\pm 1}$ or $y_j^{\pm 1}$ in the subsequence $T'_m$ of $\varphi^{(k+1)}(T'_m)$.
Of these letters of $T'_m$ which do not belong to $c$, let $t$ denote the one
furthest to the right. The folding $\C F$ must pair $t$ with some letter $s$
of $\varphi^{(k+1)}(\varphi_j(w))$. Since $s=t^{-1}$ and since the letters
$x_j^{\pm 1}$ and $y_j^{\pm 1}$ in $\varphi^{(k+1)}(\varphi_j(w))$ only occur
in the subsequences $T'_i$, it follows that $s$ is a letter in $T_i'$ for
some $1\leq i\leq m$. 
Also $s$ must appear to the left of $t$ by the choice of $t$.

Let $v_0$ be the subword of $\varphi_{j}(w)$ between the letters 
$s$ and $t$, inclusive.
Let $v_L$ be the subword of $\varphi_{j}(w)$ consisting of all letters to the left of
$s$ and $v_R$, letters to the right of~$t$, so that $v_Lv_0v_R=\varphi_{j}(w)$.
Applying Lemma \ref{L:fik-anti-triangle} to $\varphi_j(w)$ and $\varphi^{(k+1)}(\varphi_j(w))$
and $(c,\C F)$ we obtain elements $w_L=\rho_{\infty}(v_L)$, $w_0=\rho_{\infty}(v_0)$ and
$w_R=\rho_{\infty}(v_R)$ in $H = \langle x_0,x_1 \rangle$ such that
\begin{equation}
\|\varphi^{(k+1)}(\varphi_{j}(w))\|
\geq \|\varphi^{(k+1)}(w_0)\|+ \|\varphi^{(k+1)}(w_Lw_R)\|.
\label{Eq:lemma-applied}
\end{equation}
Recall from \eqref{eqn:phi_k(T_j)} that $\rho_{\infty}(T_j) = 1$.
%\rho_{\infty}\left( (x_jy_j)^n(y_jx_j)^{-n} \right)
%= (x_{j\,{\rm mod}\,2}x_{j\,{\rm mod}\,2})^n(x_{j\,{\rm mod}\,2}x_{j\,{\rm mod}\,2}%)^{-n}=1$.
Then
\begin{align*}
w_Lw_0w_R &= \rho_{\infty}(v_Lv_0v_R)\\ 
&= \rho_{\infty}(\varphi_{j}(w))\\
&=\rho_{\infty}(t_1 T'_1 t_2 T'_2\cdots t_mT'_{m})\\
&=\rho_{\infty}(t_1\cdots t_m)\\
&=t_1\cdots t_m=w
\end{align*}
It follows from
\eqref{Eq:j-notation} and\eqref{Eq:lemma-applied} that
$$
\|\varphi^{(k-1)}(w)\| \geq
\|\varphi^{(k+1)}(\varphi_{j}(w))\|
\geq \|\varphi^{(k+1)}(w_0)\|+ \|\varphi^{(k+1)}(w_Lw_R)\|.
$$
It remains to prove that $w_0\neq 1$ and $w_Lw_R\neq 1$.
There are two cases depending on the position of the letter $s$ in
$\varphi_{j}(w) = t_1T_1' \cdots t_mT_j$.

\subsubsection*{Case 2.1: {\it The letter $s$ is in the last factor $T'_m$.}}

Recall that the reduced word $T'_m=T_j = (x_jy_j)^n(x_j^{-1}y_j^{-1})^n$
is of even length. Since $s=t^{-1}$ we
get that the subword $v_0$ of $T_j$ is of odd length and therefore so is 
$w_0=\rho_{\infty}(v_0)$ because $\rho_{\infty}$ maps generators to 
generators. 
It follows that $w_0 \neq 1$.

Let $T_j- v_0$ denote the word obtained by removing $v_0$ from $T_j$.
Then
$\rho_{\infty}(T_j- v_0) = x_{j\,{\rm mod}\,2}^p$, where $p$ is an odd integer. 
It follows that
$$
w_Lw_R = \rho_{\infty}(v_Lv_R) = \rho_{\infty}(t_1T_1'\cdots t_m \cdot (T_j-v_0))
=t_1\cdots t_m x_{j\,{\rm mod}\,2}^p.
$$
Since $t_1\cdots t_m$ is reduced and $t_m=x_{(j+1)\,{\rm mod}\,2}$, the word 
$w_Lw_R=t_1\cdots t_m x_{j\,{\rm mod}\,2}^p$ is also reduced and non-trivial,
hence $w_Lw_R\neq 1$ in $H$.

\subsubsection*{Case 2.2: {\it The letter $s$ is in $T'_i$, where $1\leq i\leq m-1$.}}
Write $T'_i=T^L_i\, s\, T^R_i$ and $T'_m=T^L_m\, t\, T^R_m$.
Then
\begin{align*}
\varphi_{j}(w) 
&= 
t_1T_1'\cdots t_i T_i't_{i+1} %\underbrace{T_{i}'}_{\text{contains }s}
\cdots t_m  T_m'  %\underbrace{T'_{m}}_{\text{contains } t}\\
\\
&= 
t_1T_1'\cdots t_iT^L_i\, s\, T^R_i t_{i+1} \cdots t_m T^L_m\, t\, T^R_m.
\end{align*}
By construction 
\begin{align*}
v_L &= t_1T'_1\cdots t_iT^L_i\\
v_0 &= sT^R_it_{i+1}\cdots t_m T^L_mt\\
v_R &= T^R_m.
\end{align*}
Since $\rho_\infty(T_j)=1$ by \eqref{eqn:phi_k(T_j)}, it follows that
\begin{align*}
w_Lw_R &= \rho_{\infty}(v_Lv_R) =\rho_{\infty}(t_1T'_1\cdots t_i T^L_iT^R_m)\\
&=t_1\cdots t_i  \cdot \rho_{\infty}(T^L_iT^R_m) \\
&=t_1\cdots t_i \cdot x_{j\,{\rm mod}\,2}^p.
\end{align*}
Then $w_Lw_R$ is conjugate to $x_{j\,{\rm mod}\,2}^p \cdot t_1\cdots t_i$.
The latter word is reduced since $t_1=x_{(j+1)\,{\rm mod}\,2}$, and is  non-trivial since $i \geq 1$.
It follows that $w_Lw_R\neq 1$ in $H$. 
Similarly,
$$
w_0 = 
\rho_{\infty}(s\, T^R_it_{i+1}\cdots T_{m-1}t_mT^L_m\, t) 
= \rho_\infty(s) \cdot x_{j\,{\rm mod}\,2}^pt_{i+1}\cdots t_m x_{j\,{\rm mod}\,2}^q \cdot \rho_\infty(t)
$$
for some $p,q\in \mathbb Z$. 
Since $t=s^{-1}$ this word is conjugate to
$t_{i+1}\cdots t_m x_{j\,{\rm mod}\,2}^{p+q}$ which is reduced and non-trivial
since $i+1 \leq m$ and $t_m = x_{(j+1)\,{\rm mod}\,2}$. This implies that $w_0\neq 1$ and finishes
the proof of the lemma.
\end{proof}

\end{document}
